% ======================================================================
% capitulos/ce03_infraestructura.tex
% ----------------------------------------------------------------------
% LÍNEA CE03 — INFRAESTRUCTURA TECNOLÓGICA
% Esqueleto / marcador de posición. Reemplazar por los capítulos reales
% conforme se desarrollen los entregables de esta línea.
% ======================================================================
\chapter{Infraestructura Tecnológica — Presentación de la línea}
\label{chap:ce03-intro}

Esta línea (CE03) define la infraestructura de red y de despliegue que soporta el
sistema: la red de fibra óptica multi-zona de WISP, la interconexión de las 13 zonas
y el entorno donde se ejecuta la solución de software (línea CE02). Se apoya en la
\hyperref[chap:soluciontic]{solución TIC} definida en la línea CE01.

\begin{upeubox}{Entregables planificados de la línea CE03 (en desarrollo)}
El alcance previsto para esta línea comprende:
\begin{itemize}[leftmargin=1.4em]
  \item \textbf{Diseño de la topología de red:} arquitectura multi-zona (OLT, ONU,
  NAP, enlaces de fibra) y diagrama físico y lógico.
  \item \textbf{Direccionamiento y segmentación:} plan de direccionamiento IP, VLAN
  y separación de tráfico por zona/servicio.
  \item \textbf{Equipamiento y configuración:} routers MikroTik RB4011 (RouterOS~7),
  OLT/ONU y parámetros PPPoE.
  \item \textbf{Interconexión y acceso remoto:} VPN WireGuard para gestión
  centralizada de los 13 MikroTik desde el sistema.
  \item \textbf{Seguridad de red:} \textit{hardening}, control de accesos, firewall y
  buenas prácticas.
  \item \textbf{Monitoreo y gestión:} disponibilidad, SNMP y tableros de control de la
  red.
\end{itemize}
\end{upeubox}

\medskip
\textit{Nota: los capítulos de esta línea se incorporarán aquí a medida que se
completen sus entregables. Este archivo es un marcador de posición para mantener
integrada la estructura de las tres líneas del proyecto.}

% Al desarrollar, sustituir este archivo (o añadir \input adicionales en main.tex)
% por capítulos como:
%   \input{capitulos/ce03_topologia}
%   \input{capitulos/ce03_direccionamiento}
%   \input{capitulos/ce03_equipamiento}
%   \input{capitulos/ce03_vpn_interconexion}
%   \input{capitulos/ce03_seguridad}
%   \input{capitulos/ce03_monitoreo}
